\documentclass[../NDCNNAI_main.tex]{subfiles}
\graphicspath{{\subfix{../fig/}}}

\begin{document}
\subsection{Глубокие распределённые свёрточные сети}
В стандартном классе глубоких CNN (DCNN) по Чжоу \cite{Zhou2018} вектор смещения \emph{структурирован}. Это ключевое ограничение, которое впоследствии будет ослаблено.

\begin{definition}[Структурированный вектор смещения в $\mathcal{H}_{J}^{CNN}$ \cite{Zhou2018}]
Пусть $2 \leq s \leq d$ и $d_j = d + j s$. Для слоёв $j = 1, \ldots, J-1$ вектор смещения $b^{(j)} \in \mathbb{R}^{d_j}$ ограничен шаблоном \emph{``свободные границы -- разделяемая внутренность''} (\emph{boundary-free / interior-shared}):
\[
b^{(j)} = \left( b_1^{(j)}, \ldots, b_s^{(j)},\;
\underbrace{b_{s+1}^{(j)}, \ldots, b_{s+1}^{(j)}}_{d_j - 2s \text{ (общая часть)}},\;
b_{d_j - s+1}^{(j)}, \ldots, b_{d_j}^{(j)} \right)^T.
\]
\end{definition}

\textbf{Интерпретация:} За исключением эффектов, возникающих из-за дополнения границ (\emph{padding}), порог активации ReLU является \emph{разделяемым} (\emph{shared}) для всех пространственных позиций внутри слоя. Один и тот же скаляр $b_{s+1}^{(j)}$ используется в качестве смещения для множества нейронов. Такая структура резко ограничивает число свободных параметров смещения до $(s+1)$ на слой, но и ограничивает выразительную способность сети.

Следующим логическим шагом в работе Чжоу \cite{Zhou_Distributed} является ослабление этих ограничений на смещения.

\textbf{Ключевая мысль:} В слое со свёрточными весами (имеющими структуру матрицы Тёплица/ленточной матрицы) сеть высоко структурирована. Основной ``рычаг'', который позволяет нарушить жёсткую структуру трансляционной симметрии и создать множество различных порогов активации (``узлов'' --- \emph{knots}) в различных пространственных позициях, --- это \emph{шаблон смещений}.

Таким образом, переход от DCNN к DDCNN заключается в замене \emph{структурированного вектора} смещений на \emph{распределённую матрицу} смещений, где каждый вход связи может иметь свой собственный порог, зависящий от позиции.

В формулировке DDCNN по Чжоу \cite{Zhou_Distributed} ``смещение'' на уровне $j$ становится не вектором, а \emph{ленточной матрицей} (\emph{band matrix}), что соответствует \emph{позиционно-зависимым порогам}.

\begin{definition}[Смещение как ленточная матрица в DDCNN \cite{Zhou_Distributed}]
На уровне $j$ смещение задаётся матрицей
\[
b^{(j-1)} = \left( b_{k,i}^{(j-1)} \right) \in \mathbb{R}^{d_{j-1} \times d_j},
\]
где $b_{k,i}^{(j-1)} = 0$ вне локального окна $\max\{1, i - s\} \leq k \leq \min\{d_{j-1}, i\}$.
\end{definition}

Каждая выходная координата $i$ использует свой собственный локальный столбец смещений $(b^{(j-1)})_{.,i}$:
\[
\phi_i^{(j)}(x) = \sum_{k=\max\{1, i-s\}}^{\min\{d_{j-1}, i\}} w_{i-k}^{(j)} \sigma\left( \phi_k^{(j-1)}(x) - b_{k,i}^{(j-1)} \right), \quad 1 \leq i \leq d_j.
\]

\textbf{Связь с DCNN как частным случаем}. Если каждая строка матрицы $b^{(j-1)}$ постоянна вдоль своей ненулевой полосы, т.е. $b_{k,i}^{(j-1)} = b_{k}^{(j-1)}$ не зависит от $i$, то все столбцы $(b^{(j-1)})_{.,i}$ совпадают на активных индексах. Это немедленно даёт обычное правило обновления DCNN с единым вектором смещений $b^{(j-1)} \in \mathbb{R}^{d_{j-1}}$. Таким образом, DCNN --- это специальный случай DDCNN, а отказ от ограничения на постоянство строк даёт распределённые (позиционно-зависимые) пороги при сохранении того же самого разделения свёрточных весов.

\begin{definition}[Пространство гипотез DDCNN \cite{Zhou_Distributed}]
Фиксируем входную размерность $d \in \mathbb{N}$ и длину фильтра $s+1$, где $2 \leq s \leq d$. Ширины слоёв определяются как $d_0 = d$, $d_j = d + j s$ для $j = 1, \ldots, J$. Глубокая распределённая CNN (DDCNN) строится из следующих компонентов:
\begin{itemize}
    \item \textbf{Разделяемые локальные маски (фильтры):} $w^{(j)} = (w_0^{(j)}, \ldots, w_s^{(j)}) \in \mathbb{R}^{s+1}$ (одна и та же маска повторно используется на всех пространственных позициях).
    \item \textbf{Распределённые (позиционно-/ребро-зависимые) смещения:} $b^{(j-1)} = (b_{k,i}^{(j-1)}) \in \mathbb{R}^{d_{j-1} \times d_j}$ для $j = 2, \ldots, J$ (и финальный вектор смещений $b^{(J)} \in \mathbb{R}^{d_J}$).
    \item \textbf{Функция активации ReLU:} $\sigma(t) = t_+ = \max\{t, 0\}$.
\end{itemize}
Скрытые признаки $\{\phi_i^{(j)}(x)\}$ генерируются локальной рекурсией:
\begin{align*}
    \phi_i^{(1)}(x) &= \sum_{k=1}^{d} w_{i-k}^{(1)} x_k, \\
    \phi_i^{(j)}(x) &= \sum_{k=\max\{1, i-s\}}^{\min\{d_{j-1}, i\}} w_{i-k}^{(j)} \sigma\left( \phi_k^{(j-1)}(x) - b_{k,i}^{(j-1)} \right),
\end{align*}
для $j = 2, \ldots, J$ и $i = 1, \ldots, d_j$ (где $w_{\ell}^{(j)} = 0$ вне множества $\{0, \ldots, s\}$). Выходной слой является линейной комбинацией сдвинутых ReLU:
\[
f(x) = \sum_{k=1}^{d_J} c_k \sigma\left( \phi_k^{(J)}(x) - b_k^{(J)} \right), \quad c \in \mathbb{R}^{d_J}.
\]
Полученное пространство гипотез обозначается $\mathcal{H}_{J}^{\mathbf{w},\mathbf{b}}$.
\end{definition}

\begin{theorem}[Универсальность DDCNN, теорема 1 \cite{Zhou_Distributed}]
\label{thm:zhou_universality}
Для любого компакта $\Omega \subset \mathbb{R}^d$ и любой функции $f \in C(\Omega)$,
\[
\forall \varepsilon > 0 \; \exists J \in \mathbb{N} \; \exists (\mathbf{w}, \mathbf{b}) \; \exists g \in \mathcal{H}_{J}^{\mathbf{w},\mathbf{b}} \quad \text{такие, что} \quad \|f - g\|_{C(\Omega)} < \varepsilon.
\]
Эквивалентно, объединение пространств гипотез DDCNN плотно в $C(\Omega)$ при $J \to \infty$.
\end{theorem}

\textbf{Интерпретация (что важно запомнить):}
\begin{itemize}
    \item \textbf{Отсутствие промежуточного полносвязного перемешивания:} только локальные свёртки + ReLU, затем финальный линейный выходной слой (``readout'').
    \item \textbf{``Главный регулятор'' --- глубина:} ширины слоёв ограничены ($d_j = d + j s$ растёт только линейно).
    \item \textbf{Откуда берётся универсальность?} Из композиции локальных свёрток плюс возможность размещать множество независимых порогов ReLU через распределённые смещения.
\end{itemize}

Стандартная стратегия доказательства теорем об универсальной аппроксимации включает два шага:
\begin{enumerate}
    \item Используя классическую \textbf{теорему Стоуна-Вейерштрасса}, приблизить целевую непрерывную функцию $f \in C(\Omega)$ многочленом $P$:
    \[
    \sup_{x \in \Omega} |f(x) - P(x)| \leq \varepsilon/2.
    \]
    \item Показать, что DDCNN могут равномерно приблизить $P$ на $\Omega$ с точностью $\varepsilon/2$.
\end{enumerate}

Таким образом, техническим ядром доказательства становится вопрос: \textbf{``Может ли DDCNN аппроксимировать многомерные полиномы?''} Положительный ответ на этот вопрос, основанный на конструировании кусочно-линейных сплайнов с использованием распределённых смещений, и завершает доказательство теоремы \ref{thm:zhou_universality}.

Для доказательства теоремы \ref{thm:zhou_universality} необходимо показать, как DDCNN может аппроксимировать произвольный многомерный полином. Это достигается в три основных этапа, опирающихся на следующие ключевые леммы из статьи \cite{Zhou_Distributed}.

\begin{lemma}[Разложение полинома на одномерные \cite{Zhou_Distributed}]
\label{lem:poly_decomp}
Пусть $d, \Gamma \in \mathbb{N}$. Существует множество векторов $\{\xi_k\}_{k=1}^{n_\Gamma} \subset \mathbb{R}^d$ с $\|\xi_k\|_2 = 1$, где $n_\Gamma = \binom{d-1+\Gamma}{\Gamma}$, такое что любой полином $P_\Gamma \in \mathcal{P}_\Gamma(\mathbb{R}^d)$ степени не выше $\Gamma$ можно представить в виде
\[
P_\Gamma(x) = \sum_{k=1}^{n_\Gamma} p_{k,\Gamma}(\xi_k \cdot x), \quad x \in \mathbb{R}^d,
\]
где каждое $p_{k,\Gamma} \in \mathcal{P}_\Gamma(\mathbb{R})$ является одномерным полиномом.
\end{lemma}

\begin{lemma}[Аппроксимация сплайнами \cite{Zhou_Distributed}]
\label{lem:spline_approx}
Пусть $t = \{t_1 < t_2 < \cdots < t_{N-1} < t_N\}$ --- последовательность узлов. Определим линейный оператор $L_t$ на $C[t_2, t_{N-1}]$:
\[
L_t(f)(u) = \sum_{j=2}^{N-1} f(t_j) \delta_j(u), \quad u \in [t_2, t_{N-1}],
\]
где
\[
\delta_j(u) = \frac{1}{t_j - t_{j-1}} \sigma(u - t_{j-1}) - \frac{t_{j+1} - t_{j-1}}{(t_{j+1} - t_j)(t_j - t_{j-1})} \sigma(u - t_j) + \frac{1}{t_{j+1} - t_j} \sigma(u - t_{j+1}).
\]
Тогда для любой $f \in C[t_2, t_{N-1}]$ справедлива оценка
\[
\|L_t(f) - f\|_{C[t_2, t_{N-1}]} \leq 2 \omega(f, \Delta_t),
\]
где $\Delta_t = \max_{j=3,\ldots,N-1} \{|t_j - t_{j-1}|\}$, а $\omega(f, \mu)$ --- модуль непрерывности функции $f$.
\end{lemma}

\begin{lemma}[Факторизация фильтров \cite{Zhou_Distributed}]
\label{lem:filter_factorization}
Пусть $s \geq 2$ и $W: \mathbb{Z} \to \mathbb{R}$ --- маска фильтра с носителем на $\{0, 1, \ldots, S_W\}$, причём $W_0 > 0$ и $W_{S_W} \neq 0$, где $S_W \geq d \geq s$. Тогда существует целое число $I \in [\frac{S_W}{s}, \frac{S_W}{s-1} + 2)$ и последовательность масок фильтров $\{w^{(j)}\}_{j=1}^I$ с носителем на $\{0, 1, \ldots, s\}$, где $w_0^{(1)} = W_0 > 0$ и $w_0^{(j)} = 1$ для $j = 2, \ldots, I$, таких что выполняется свёрточная факторизация:
\[
W = w^{(I)} * w^{(I-1)} * \cdots * w^{(2)} * w^{(1)}.
\]
\end{lemma}

\begin{proof}
    Пусть $\epsilon > 0$. Тогда существует полином $P_\Gamma \in \mathcal{P}_\Gamma(\mathbb{R}^d)$ с некоторым $\Gamma \in \mathbb{N}$ такой, что
    \[
    \|f - P_\Gamma\|_{C(\Omega)} \leq \frac{\epsilon}{2}.
    \]

    По лемме \ref{lem:poly_decomp} (разложение полинома на одномерные) существует множество векторов $\{\xi_k\}_{k=1}^{n_\Gamma} \subset \mathbb{R}^d$ такое, что $\|\xi_k\|_2 = 1$ для каждого $k$ и полином $P_\Gamma \in \mathcal{P}_\Gamma(\mathbb{R}^d)$ может быть выражен в виде
    \[
    P_\Gamma(x) = \sum_{k=1}^{n_\Gamma} p_{k,\Gamma}(\xi_k \cdot x), \quad x \in \mathbb{R}^d,
    \]
    где $\{p_{k,\Gamma}\}_{k=1}^{n_\Gamma} \subset \mathcal{P}_\Gamma(\mathbb{R})$ являются одномерными полиномами.

    Заметим, что $|\xi_k \cdot x| \leq \|\xi_k\|_2 \|x\|_2 \leq \|x\|_2$ для каждого $k \in \{1, \ldots, n_\Gamma\}$. Следовательно, для $x \in \Omega$ имеем $|\xi_k \cdot x| \leq B_2^{(0)}$, где
    \[
    B_2^{(0)} \coloneq \max_{x \in \Omega} \|x\|_2 < \infty.
    \]

    Таким образом, мы рассматриваем аппроксимацию одномерных функций $\{p_{k,\Gamma}\}$ на отрезке $[-B_2^{(0)}, B_2^{(0)}]$. Для $4 \leq N \in \mathbb{N}$ возьмем последовательность узлов $t = \{t_1 < t_2 < \cdots < t_{N-1} < t_N\}$ как
    \[
    t_j = -B_2^{(0)} + (j - 2)\frac{2B_2^{(0)}}{N - 3}, \quad j = 1, 2, \ldots, N-1, N,
    \]
    что влечет $[t_2, t_{N-1}] = [-B_2^{(0)}, B_2^{(0)}]$ и $\Delta_t = \frac{2B_2^{(0)}}{N - 3}$.

    По лемме \ref{lem:spline_approx} (аппроксимация сплайнами) существует линейный оператор $L_t$ такой, что
    \[
    \|L_t(p_{k,\Gamma}) - p_{k,\Gamma}\|_{C[-B_2^{(0)}, B_2^{(0)}]} \leq 2\omega\left( p_{k,\Gamma}, \frac{2B_2^{(0)}}{N - 3} \right),
    \]
    где $\omega(f, \mu)$ --- модуль непрерывности функции $f$.

    Поскольку $\lim_{\mu \to 0^+} \omega(g, \mu) = 0$ для любой $g \in C[-B_2^{(0)}, B_2^{(0)}]$, мы знаем, что существует некоторое $\mu_\epsilon > 0$ такое, что $\omega(p_{k,\Gamma}, \mu) \leq \frac{\epsilon}{4n_\Gamma}$ для любого $0 < \mu \leq \mu_\epsilon$ и $k = 1, \ldots, n_\Gamma$. Возьмем некоторое $N \in \mathbb{N}$ такое, что $\frac{2B_2^{(0)}}{N - 3} \leq \mu_\epsilon$ и $N \geq 4$. Тогда мы имеем
    \[
    \left\| \sum_{k=1}^{n_\Gamma} L_t(p_{k,\Gamma})(\xi_k \cdot x) - P_\Gamma \right\|_{C(\Omega)} \leq \sum_{k=1}^{n_\Gamma} \|L_t(p_{k,\Gamma})(u) - p_{k,\Gamma}(u)\|_{C[-B_2^{(0)}, B_2^{(0)}]} \leq \frac{\epsilon}{2}.
    \]

    Теперь обратимся к построению маски фильтра $W$. Выразим каждый вектор $\xi_k$ через его компоненты
    \[
    \xi_k = ((\xi_k)_\ell)_{\ell=1}^d \in \mathbb{R}^d, \quad k = 1, \ldots, n_\Gamma.
    \]

    Определим маску фильтра $W : \mathbb{Z} \to \mathbb{R}$ с носителем на $\{0, 1, \ldots, dn_\Gamma + 1\}$, положив $W_0 = W_{dn_\Gamma + 1} = 1$ и
    \[
    W_{(k-1)d+\ell} = (\xi_k)_{d+1-\ell}, \quad \ell = 1, \ldots, d, \quad k = 1, \ldots, n_\Gamma.
    \]

    Затем мы применяем лемму \ref{lem:filter_factorization} (факторизация фильтров) к маске фильтра $W : \mathbb{Z} \to \mathbb{R}$ с $S_W = dn_\Gamma + 1$. Существует целое число $I \in [\frac{dn_\Gamma + 1}{s}, \frac{dn_\Gamma + 1}{s-1} + 2)$ и последовательность масок фильтров $\{w^{(j)}\}_{j=1}^I$ с носителем на $\{0, 1, \ldots, s\}$, имеющая $w_0^{(1)} = W_0 = 1$ и $w_0^{(j)} = 1$ для $j = 2, \ldots, I$, таких что выполняется свёрточная факторизация:
    \[
    W = w^{(I)} * w^{(I-1)} * \cdots * w^{(2)} * w^{(1)}.
    \]

    Комбинируя это с построением инициализирующих слоёв DDCNN (раздел 3 статьи), мы получаем, что после $I$ слоёв компонента с индексом $kd+1$ гомогенизированного представления равна
    \[
    \hat{\phi}_{kd+1}^{(I)}(x) = \xi_k \cdot x, \quad k = 1, \ldots, n_\Gamma.
    \]

    Используя построение глубоких слоёв DDCNN (раздел 4 статьи) с дополнительными $N-1$ группами по $I$ слоёв (всего $J = NI$ слоёв) и специальным выбором масок $\{w^{(j)}\}_{j=I+1}^{NI}$ и матриц смещений, мы получаем, что в последнем слое $\Phi^{(J)}(x)$ содержатся функции вида
    \[
    \sigma\left(\hat{\phi}_{d+\ell}^{(I)}(x) - t_i\right), \quad \ell = 1, \ldots, Js, \quad i = 1, \ldots, N.
    \]

    В частности, взяв $\ell = (k-1)d + 1$ для $k = 1, \ldots, n_\Gamma$, мы получаем, что
    \[
    \sigma(\xi_k \cdot x - t_i) \in \operatorname{span} \left\{ \Phi_i^{(J)}(x) \right\}_{i=d+1}^{d+Js} \subseteq \mathcal{H}_J^{\mathbf{w},\mathbf{b}}.
    \]

    Но определение линейного оператора $L_t$ (лемма \ref{lem:spline_approx}) говорит нам, что
    \[
    L_t(p_{k,\Gamma})(\xi_k \cdot x) \in \operatorname{span} \{\sigma(\xi_k \cdot x - t_i), i = 1, \ldots, N\}
    \]
    для каждого $k \in \{1, \ldots, n_\Gamma\}$. Следовательно,
    \[
    \sum_{k=1}^{n_\Gamma} L_t(p_{k,\Gamma})(\xi_k \cdot x) \in \operatorname{span} \{\sigma(\xi_k \cdot x - t_i), k = 1, \ldots, n_\Gamma, i = 1, \ldots, N\} \subseteq \mathcal{H}_J^{\mathbf{w},\mathbf{b}}.
    \]

    Таким образом, взяв $f^* = \sum_{k=1}^{n_\Gamma} L_t(p_{k,\Gamma})(\xi_k \cdot x) \in \mathcal{H}_J^{\mathbf{w},\mathbf{b}}$, мы получаем
    \[
    \|f - f^*\|_{C(\Omega)} \leq \|f - P_\Gamma\|_{C(\Omega)} + \|P_\Gamma - f^*\|_{C(\Omega)} \leq \frac{\epsilon}{2} + \frac{\epsilon}{2} = \epsilon.
    \]

    Это доказывает, что для любого $\epsilon > 0$ существует глубина $J = NI$ и параметры $(\mathbf{w}, \mathbf{b})$ такие, что $\mathcal{H}_J^{\mathbf{w},\mathbf{b}}$ содержит функцию $f^*$, аппроксимирующую $f$ с точностью $\epsilon$, что и утверждает теорема \ref{thm:zhou_universality}.
\end{proof}


\textbf{Ключевые моменты доказательства:}
\begin{enumerate}
    \item Аппроксимация непрерывной функции полиномом (теорема Стоуна–Вейерштрасса).
    \item Разложение многомерного полинома на сумму одномерных полиномов от скалярных произведений $\xi_k \cdot x$ (лемма \ref{lem:poly_decomp}).
    \item Аппроксимация одномерных полиномов кусочно-линейными сплайнами $L_t(p_{k,\Gamma})$, построенными из сдвигов ReLU $\sigma(u - t_i)$ (лемма \ref{lem:spline_approx}).
    \item Кодирование направлений $\xi_k$ в маску фильтра $W$ и её факторизация на короткие фильтры $w^{(j)}$ длины $s+1$ (лемма \ref{lem:filter_factorization}).
    \item Использование инициализирующих слоёв DDCNN для генерации линейных функций $\xi_k \cdot x$.
    \item Использование глубоких слоёв DDCNN с распределёнными смещениями для независимого размещения порогов $t_i$ и генерации функций $\sigma(\xi_k \cdot x - t_i)$.
    \item Линейная комбинация этих функций даёт сплайн $L_t(p_{k,\Gamma})(\xi_k \cdot x)$, который принадлежит $\mathcal{H}_J^{\mathbf{w},\mathbf{b}}$.
\end{enumerate}

Далее проясним, какая интуитивная идея, стоит за переходом от стандартных DCNN к распределённым DDCNN. Мы рассмотрим, как механизм распределённых (позиционно-зависимых) смещений позволяет преодолеть фундаментальные ограничения классической свёртки и создавать сложные функции аппроксимации.

Чистая свёрточная операция (без смещений) обладает свойством \emph{разделения весов} (\emph{weight sharing}): один и тот же локальный линейный функционал применяется в каждой позиции.
Это создаёт сильные структурные ограничения. В стандартном DCNN с разделяемыми смещениями (\emph{shared bias}) ситуация несколько улучшается, но остаётся фундаментальная проблема:
\begin{itemize}
    \item \textbf{Разделяемые смещения:} $b_{k,i}^{(j-1)} \equiv \beta_k^{(j-1)}$ (нет зависимости от целевой позиции $i$)
    \item \textbf{Следствие:} функция активации $\sigma\left( \phi_k^{(j-1)}(x) - b_{k,i}^{(j-1)} \right)$ использует один и тот же порог везде, где повторно используется $\phi_k^{(j-1)}$.
    \item \textbf{Проблема:} точки излома (\emph{breakpoints}) и паттерны активации глобально выровнены по позициям (с точностью до сдвигов), что ограничивает то, насколько мелко мы можем ``размещать узлы'' при фиксированных локальных масках.
\end{itemize}

\textbf{Ключевой вопрос:} можем ли мы выбирать различные положения ``узлов'' ReLU (порогов) в разных позициях $i$?
В DDCNN ослабляется ограничение на разделение смещений:
\begin{itemize}
    \item \textbf{Распределённые смещения:} $b_{k,i}^{(j-1)}$ свободны (в пределах локальной полосы $i-s \leq k \leq i$)
    \item \textbf{Следствие:} каждое соединение, ведущее в позицию $i$, может иметь свой собственный порог:
    \[
    \phi_i^{(j)}(x) = \sum_{k=\max\{1, i-s\}}^{\min\{d_{j-1}, i\}} w_{i-k}^{(j)} \sigma\left( \phi_k^{(j-1)}(x) - b_{k,i}^{(j-1)} \right).
    \]
\end{itemize}
\textbf{Что это даёт:} мы можем активировать/подавлять один и тот же разделяемый признак $\phi_k^{(j-1)}$ по-разному в разных пространственных позициях $i$. Это именно то, что нужно для построения множества локализованных паттернов ``узлов'' и (с помощью глубины) сборки богатого базиса кусочно-линейных функций.

\textbf{Интерпретация (1D случай):} Рассмотрим $b_{k,i}^{(j-1)}$ как \emph{расписание порогов, зависящее от позиции}. Даже если маска $w^{(j)}$ разделяемая, сеть может решать, \emph{где} (в какой позиции $i$) ReLU ``включается''. Глубина затем собирает эти локально управляемые отклики в произвольную непрерывную целевую функцию.

\begin{remark}
    DDCNN = разделяемые локальные фильтры $w^{(j)}$ \textbf{какой признак} + распределённые пороги $b^{(j)}$ \textbf{где он активируется}.
\end{remark}

\textbf{Интерпретация:}
\begin{itemize}
    \item \textbf{Разделяемые фильтры ($w^{(j)}$):} определяют, \emph{какие локальные признаки} извлекаются (например, границы, текстуры).
    \item \textbf{Распределённые смещения ($b^{(j)}$):} определяют, \emph{в каких пространственных позициях} эти признаки становятся релевантными (активируются).
    \item \textbf{Глубина:} позволяет комбинировать эти локально-управляемые активации для создания всё более сложных и глобальных зависимостей.
\end{itemize}

ReLU-сети естественным образом подходят для \emph{кусочно-линейных} конструкций. Это используется в доказательстве универсальности Чжоу.
\textbf{Ключевым примитивом} является представление/аппроксимация функций вида
\[
x \mapsto (x - t)_+, \quad \text{где } (x)_+ \coloneq \max\{x, 0\}.
\]
Из этих элементарных ``строительных блоков'' можно конструировать:
\begin{itemize}
    \item \textbf{Тент-функции / шляпки} (\emph{tent functions / hat functions})
    \item \textbf{Линейные сплайны} (\emph{linear splines})
    \item \textbf{B-сплайноподобные базисы} (через линейные комбинации и композиции)
    \item \textbf{Аппроксимации полиномов} с использованием сплайн-базисов
\end{itemize}

\begin{remark}
    В конструкции Чжоу \emph{свёрточная структура} эффективно представляет сдвинутые копии, в то время как \emph{распределённые смещения} задают узлы (\emph{knots}).
\end{remark}
\textbf{Замечание:} 

Как одномерная свёрточная архитектура может выражать многомерные зависимости?
В доказательствах аппроксимации CNN используются две распространённые математические техники:

\begin{itemize}
    \item \textbf{Кодирование индексами:} рассматривать входной вектор $x \in \mathbb{R}^d$ как ``последовательность''; свёртка агрегирует локальные группы, а глубина увеличивает область взаимодействия (рост рецептивного поля).
    
    \item \textbf{Композиционные строительные блоки:} полиномы могут быть выражены через композиции сложений и умножений, которые могут быть аппроксимированы кусочно-линейными сетями.
\end{itemize}

\textbf{Вывод:} глубина $J$ создаёт постепенно более глобальные взаимодействия, в то время как фильтры остаются маленькими.

\subsection{Распределённые свёрточные сети и операторное обучение}
Теперь установим мост между теорией универсальной аппроксимации глубоких свёрточных сетей и практической задачей обучения операторов. Мы рассмотрим, когда свёрточная архитектура является подходящим индуктивным смещением для операторного обучения, и сравним её с альтернативными подходами.

С точки зрения обучения операторов мы изучаем \textbf{оператор} (отображение между функциональными пространствами):
\[
G\colon X \to Y, \quad X = C(D; \mathbb{R}^{d_{\text{in}}}) \text{ или } L^2(D; \mathbb{R}^{d_{\text{in}}}), \quad Y = C(U; \mathbb{R}^{d_{\text{out}}}) \text{ или } L^2(U; \mathbb{R}^{d_{\text{out}}}),
\]
по парам $\{(u^{(n)}, G(u^{(n)}))\}_{n=1}^N$.

В DeepONet (лекции 4--6) гипотеза представлялась как $N = R \circ A \circ E$:
\begin{itemize}
    \item \textbf{Энкодер / сенсоры} $E : X \to \mathbb{R}^m$ (или $\mathbb{R}^{md_{\text{in}}}$):
    \[
    E(u) = (u(x_1), \ldots, u(x_m)) \quad \text{(точечные сенсоры)}
    \]
    или
    \[
    E(u) = (\text{Patch}_r(u; y_1), \ldots) \quad \text{(локальные патчи)}.
    \]
    
    \item \textbf{Конечномерный аппроксиматор} $A : \mathbb{R}^m \to \mathbb{R}^p$ (``ветвистая сеть''):
    \[
    A(E(u)) = \alpha \in \mathbb{R}^p.
    \]
    
    \item \textbf{Реконструктор} $R : \mathbb{R}^p \to Y$ (``стволовая сеть / базис''):
    \[
    (R(\alpha))(y) = \sum_{k=1}^p \alpha_k \tau_k(y),
    \]
    где $\tau_k$ --- выученные/параметрические базисные функции (часто через стволовую сеть).
\end{itemize}

Встаёт вопрос, когда $E$ и/или $A$ должны быть \textbf{свёрточными}, а не полносвязными?
Рассмотрим оператор сдвига: $(R_\gamma u)(x) = u(x - \gamma)$. Оператор $G$ называется \emph{эквивариантным относительно сдвига}, если
\[
G(R_\gamma u) = R_\gamma G(u).
\]
Если целевой оператор приближённо соблюдает такие симметрии и является локальным, то \emph{разделение весов свёртки} может быть правильным индуктивным смещением.

Свёрточные / локальные архитектуры, как правило, помогают, когда оператор $G\colon X \to Y$ (приближённо) соответствует следующим структурным свойствам:

\begin{enumerate}
    \item \textbf{Эквивариантность относительно сдвига (поля $\to$ поля):}
    \[
    G(R_\gamma u) \approx R_\gamma G(u) \quad \text{для сдвигов } (R_\gamma u)(x) = u(x - \gamma).
    \]
    Типично для однородных операторов решения УЧП на периодической области / больших внутренних областях.
    
    \item \textbf{Инвариантность относительно сдвига (поля $\to$ скаляры/статистики):} существует функционал $\Phi : X \to \mathbb{R}$ с $\Phi(R_\gamma u) = \Phi(u)$ (например, глобальные сводки, энергии, метки классов после глобального пулинга).
    
    \item \textbf{Локальная зависимость (конечное рецептивное поле):} существует $r > 0$ такое, что, грубо говоря,
    \[
    u|_{B_r(y)} = v|_{B_r(y)} \implies (G(u))(y) \approx (G(v))(y),
    \]
    то есть выход в позиции $y$ зависит в основном от входа вблизи $y$.
    
    \item \textbf{Стационарность / разделение весов по пространству:} одно и то же ``локальное правило'' применяется во многих позициях (один и тот же ядр везде).
    
    \item \textbf{Мультимасштабная структура:} релевантные паттерны проявляются от мелкого к крупному масштабу; субдискретизация/пулинг строит иерархии.
\end{enumerate}

Примерами $G$ могут выступать операторы УЧП с постоянными коэффициентами на сетках; преобразование изображений (шумоподавление, сегментация); локальное пространственно-временное прогнозирование; замыкания турбулентности с локальными шаблонами.

Свёртка может быть ограничением, когда нарушается одно (или более) из следующих условий:

\begin{enumerate}
    \item \textbf{Отсутствие естественной группы сдвигов:} нерегулярные области (графы / многообразия / облака точек) или сильно варьирующаяся плотность дискретизации.
    
    \item \textbf{Сильная нестационарность / неоднородность:} коэффициенты или статистики меняются с позицией (например, УЧП с переменными коэффициентами), поэтому глобальное разделение весов является неверной спецификацией без позиционных признаков/эмбеддингов.
    
    \item \textbf{Глобальная (дальнодействующая) связь доминирует:} интегральные операторы с широкими ядрами, зависимости типа внимания (\emph{attention}), или эллиптические УЧП, где эффекты функции Грина индуцируют нелокальность.
    
    \item \textbf{Граничные эффекты существенны:} симметрия сдвигов нарушается вблизи границ; выбор дополнения (\emph{padding}) имеет значение; оператор зависит от граничных условий.
    
    \item \textbf{Разреженные сенсоры, а не плотные сетки:} если доступно лишь несколько измерений, глобальный базис (DeepONet/FMLP) может быть более эффективным по данным.
\end{enumerate}

\textbf{Альтернативы / дополнения:} DeepONet/FMLP с выученными глобальными базисами; фурье/спектральные операторные модели (в стиле FNO); внимание/трансформеры для дальнодействующей связи; графовые нейронные операторы для нерегулярных областей.

Сравним концепты DeepONet и свёрточных операторных моделей
\begin{table}[h]
\centering
\begin{tabular}{|p{0.2\textwidth}|p{0.35\textwidth}|p{0.35\textwidth}|}
\hline
\textbf{Аспект} & \textbf{DeepONet / FMLP} & \textbf{Свёрточная операторная модель} \\
\hline
Обработка входа & Сенсоры, сэмплы $E(u) = (u(x_i))_{i=1}^m$ & Сетка, патчи + свёрточный экстрактор признаков \\
\hline
Механизм смешивания & Глобальное смешивание в ветвистой сети $A\colon \mathbb{R}^m \to \mathbb{R}^p$ & Локальное смешивание через ядра (разделение весов), глубина строит большое рецептивное поле \\
\hline
Обработка выхода & Ствол/базис $R(\alpha)(y) = \sum_{k=1}^p \alpha_k \tau_k(y)$ & Свёрточный декодер, апсемплинг, покоординатная голова (часто эквивариантная) \\
\hline
Встроенное смещение & Глобальный низкоразмерный базис; отсутствие эквивариантности сдвигов, если не сконструирована явно & Локальность + (приближённая) эквивариантность сдвигов + разделение параметров \\
\hline
Лучше всего, когда & Гладкая глобальная структура; разреженные сенсоры; обучение, не зависящее от геометрии & Локальные, стационарные, мультимасштабные поля; плотные сетки; шаблоны УЧП, операторы типа изображений \\
\hline
Основной риск & Неверная спецификация базиса; может игнорировать локальную физику без индуктивной структуры & Пропуск глобальных взаимодействий, граничных эффектов, нестационарности без дополнительных механизмов \\
\hline
\end{tabular}
\caption{Концептуальное сравнение DeepONet/FMLP и свёрточных операторных моделей}
\label{tab:deeponet_vs_conv}
\end{table}
\textbf{Вывод:} выбирайте класс гипотез в соответствии с симметриями оператора. Можно также \emph{совмещать} их: использовать CNN внутри $E$ или $A$, сохраняя при этом стиль DeepONet для $R$.
Более подробно:
\begin{itemize}
    \item \textbf{Энкодер $E$ (как вы представляете входную функцию):}
    \begin{itemize}
        \item Сенсоры: $E(u) = (u(x_i))_{i=1}^m$ (классический DeepONet),
        \item Патчи / локальные окна: $E(u) = (\text{Patch}_r(u; y_i))_{i=1}^m$ (удобно для CNN),
        \item CNN-энкодер на сетке для получения поля признаков (эквивариантное представление).
    \end{itemize}
    
    \item \textbf{Аппроксиматор $A$ (конечное отображение):}
    \begin{itemize}
        \item MLP (глобальное смешивание) vs CNN/ResNet (локальное смешивание + глубина),
        \item опционально: добавить дальнодействующее смешивание (фурье-слой / внимание), если физика нелокальна.
    \end{itemize}
    
    \item \textbf{Реконструктор $R$ (как вы производите выход):}
    \begin{itemize}
        \item \emph{Глобальный базис} (ствол): $R(\alpha)(y) = \sum_{k=1}^p \alpha_k \pi_k(y)$,
        \item \emph{Свёрточный декодер}: стиль U-Net апсемплинг / покоординатная голова (локальные выходы).
    \end{itemize}
\end{itemize}

\begin{remark}[Правило проектирования]
    Обучение операторов не обязывает использовать полносвязные компоненты: выбирайте $E/A/R$ в соответствии с симметриями оператора (эквивариантность/инвариантность) и локальностью. Гибридные модели позволяют сохранить глобальную выразительность, одновременно навязывая локальную структуру там, где это важно.    
\end{remark}

\subsection{Основы RKHS: ядра как скалярные произведения (минимальный инструментарий)}
В данном разделе вводится математический аппарат воспроизводящих гильбертовых пространств (RKHS) --- ключевой теоретический инструмент, позволяющий связать ядерные методы с анализом свойств свёрточных нейронных сетей. Мы определим положительно определённые ядра, сформулируем фундаментальную теорему Мура-Аронсайна и объясним свойство воспроизведения --- основу для анализа устойчивости и обобщения.

\begin{definition}[Положительно определённое ядро]
\label{def:pd_kernel}
Функция \( K\colon \mathcal{X} \times \mathcal{X} \to \mathbb{R} \) называется \emph{положительно определённой}, если для любого конечного набора точек \(\{x_1, \ldots, x_n\} \subset \mathcal{X}\) и любых коэффициентов \( c = (c_1, \ldots, c_n) \in \mathbb{R}^n \) выполняется неравенство:
\[
\sum_{i,j=1}^{n} c_i c_j K(x_i, x_j) \geq 0.
\]
Эквивалентно, \emph{матрица Грама} \( \mathbf{K} \coloneq [K(x_i, x_j)]_{i,j=1}^n \) является положительно полуопределённой: \( \mathbf{K} \succeq 0 \).
\end{definition}
\textbf{Интуиция:} Положительная определённость ядра гарантирует, что его можно интерпретировать как \emph{скалярное произведение} в некотором (возможно, бесконечномерном) гильбертовом пространстве, даже если это произведение вычисляется в исходном пространстве \(\mathcal{X}\) \emph{неявно}.

\begin{theorem}[Мура-Аронсайна]
\label{thm:moore_aronszajn}
Для любого положительно определённого ядра \( K \) на множестве \(\mathcal{X}\) существует единственное гильбертово пространство функций \(\mathcal{H}_K \subset \mathbb{R}^{\mathcal{X}}\), называемое \emph{воспроизводящим гильбертовым пространством (RKHS)}, такое что:
\begin{enumerate}
    \item Для любого \( x \in \mathcal{X} \) функция \( K_x \coloneq K(x, \cdot) \) принадлежит \(\mathcal{H}_K\).
    \item Для любой функции \( f \in \mathcal{H}_K \) и любой точки \( x \in \mathcal{X} \) выполняется \emph{свойство воспроизведения} (\emph{reproducing property}):
    \[
    f(x) = \langle f, K_x \rangle_{\mathcal{H}_K}.
    \]
    \item Линейная оболочка функций \(\{K_x : x \in \mathcal{X}\}\) плотна в \(\mathcal{H}_K\).
\end{enumerate}
Пространство \(\mathcal{H}_K\) часто называют \emph{пространством признаков}, а отображение \(\Phi\colon \mathcal{X} \to \mathcal{H}_K\), заданное как \(\Phi(x) = K_x\), --- \emph{отображением в признаковое пространство} (\emph{feature map}). При этом ядро выражается через скалярное произведение: \( K(x, x') = \langle \Phi(x), \Phi(x')\rangle_{\mathcal{H}_K} \).
\end{theorem}
Эта теорема является краеугольным камнем ядерных методов, поскольку гарантирует, что любое положительно определённое ядро корректно задаёт гильбертово пространство функций с удобным свойством воспроизведения.

Свойство воспроизведения \( f(x) = \langle f, K_x \rangle_{\mathcal{H}_K} \) имеет два фундаментальных следствия:
\begin{itemize}
    \item \textbf{Непрерывность функционала вычисления:} Для любого \( x \in \mathcal{X} \) функционал вычисления \( \delta_x: f \mapsto f(x) \) является линейным и \emph{непрерывным} на \(\mathcal{H}_K\), так как \( |\delta_x(f)| = |f(x)| \leq \|f\|_{\mathcal{H}_K} \|K_x\|_{\mathcal{H}_K} \).
    \item \textbf{Ядерная машина (Kernel Machine):} Любая функция из \(\mathcal{H}_K\) может быть представлена как линейный функционал в признаковом пространстве: \( f(x) = \langle w, \Phi(x) \rangle_{\mathcal{H}_K} \) для некоторого \( w \in \mathcal{H}_K \). Это позволяет строить \emph{линейные модели в \(\mathcal{H}_K\)}, которые соответствуют сложным нелинейным решающим правилам в исходном пространстве \(\mathcal{X}\).
\end{itemize}

Аппарат RKHS предоставляет единую основу для решения задач обучения и анализа моделей.
\begin{enumerate}
    \item \textbf{Норма как мера сложности (регуляризация).} Стандартная задача обучения в RKHS формулируется как минимизация функционала риска с регуляризацией:
    \[
    \min_{f \in \mathcal{H}_K} \frac{1}{n} \sum_{i=1}^{n} \ell(y_i, f(x_i)) + \lambda \|f\|^2_{\mathcal{H}_K},
    \]
    где \( \|f\|_{\mathcal{H}_K} \) служит естественной мерой сложности (\emph{smoothness}) функции. \textbf{Теорема представления} гарантирует, что решение такой задачи лежит в конечномерной линейной оболочке ядер, центрированных на данных: \( f(\cdot) = \sum_{i=1}^n \alpha_i K(x_i, \cdot) \) \cite{Bietti_Group}.

    \item \textbf{Устойчивость ``проходит через'' признаковое отображение.} Используя неравенство Коши-Буняковского и свойство воспроизведения, можно получить оценку устойчивости предсказаний:
    \[
    |f(x) - f(x')| = |\langle f, \Phi(x) - \Phi(x')\rangle_{\mathcal{H}_K}| \leq \|f\|_{\mathcal{H}_K} \|\Phi(x) - \Phi(x')\|_{\mathcal{H}_K}.
    \]
    Таким образом, если само \emph{признаковое отображение} \(\Phi\) устойчиво (например, к малым деформациям или сдвигам) и норма \(\|f\|_{\mathcal{H}_K}\) контролируется, то и предсказания модели будут устойчивыми.
\end{enumerate}

\textbf{Мост к конструкции Бьетти--Майраль \cite{Bietti_Group}:} Если спроектировать ядро \(K\) (и, соответственно, признаковое отображение \(\Phi\)) так, чтобы оно кодировало свёрточную структуру и желаемые симметрии, то анализ инвариантности, устойчивости и обобщающей способности соответствующей модели можно проводить, используя геометрию RKHS.

\subsection{Подход Бьетти--Майраль: ядерное представление свёрточных архитектур }
\label{sec:5.2_conv_kernel_view}
Дальше излагается ключевая идея работы Бьетти и Майраля \cite{Bietti_Group}: представление глубокой свёрточной сети (CNN) как многоуровневого признакового отображения в воспроизводящее гильбертово пространство (RKHS). Этот формализм позволяет строго анализировать такие свойства CNN, как инвариантность, устойчивость и сложность, используя аппарат теории ядер.

Бьетти и Майраль предлагают строить представление \(\Phi_n\) рекуррентно, имитируя слои CNN, что порождает соответствующее ядро глубины \(n\) \cite{Bietti_Group}.
\begin{definition}[Многоуровневое признаковое отображение и ядро]
\label{def:deep_conv_kernel}
Пусть \(x_0 \coloneq x\) --- входной сигнал. Для каждого уровня \(k = 1, \ldots, n\) определяем:
\[
x_k \coloneq A_k \, M_k \, P_k \, x_{k-1},
\]
где
\begin{itemize}
    \item \(P_k\) --- оператор \emph{извлечения локальных патчей} (аналог свёрточного фильтра).
    \item \(M_k\) --- \emph{нелинейное отображение / нормализация} на уровне патча, индуцированное некоторым \emph{ядром скалярного произведения} (dot-product kernel) в пространстве признаков \(\mathcal{H}_k\).
    \item \(A_k\) --- \emph{линейный оператор усреднения (пулинга)} (например, локальное усреднение или субдискретизация).
\end{itemize}
Глубинное признаковое отображение определяется как \(\Phi_n(x) \coloneq x_n\).
Соответствующее \emph{глубинное ядро} задаётся скалярным произведением в финальном пространстве признаков:
\[
K_n(x, x') \coloneq \langle \Phi_n(x), \Phi_n(x')\rangle_{\mathcal{H}_n} = \langle x_n, x_n'\rangle_{\mathcal{H}_n}.
\]
Это ядро порождает RKHS \(\mathcal{H}_{K_n}\).
\end{definition}
Интуитивно каждый шаг \(P_k \mapsto M_k \mapsto A_k\) соответствует слою CNN: линейная свёртка (\(P_k\)), нелинейная функция активации (\(M_k\)), операция пулинга (\(A_k\)). Однако в ядерной постановке нелинейность \(M_k\) трактуется как отображение в (возможно, бесконечномерное) RKHS, связанное с ядром на патчах.

В теоретической конструкции Бьетти-Майраля \cite{Bietti_Group} пространство признаков для каждого патча является RKHS (возможно, бесконечномерным). Конечноканальная свёрточная сеть (CNN) или свёрточная ядерная сеть (CKN) рассматривается как \emph{приближение или проекция} этого бесконечномерного признакового отображения на конечное число направлений.

Это делает \emph{геометрию представления явной и доступной для анализа}: свойства всего представления \(\Phi_n\) выводятся из свойств элементарных операторов \(P_k, M_k, A_k\) и геометрии RKHS.

Важный результат \cite{Bietti_Group} заключается в том, что построенное RKHS \(\mathcal{H}_{K_n}\) содержит в себе функции, реализуемые многими классическими CNN.
\begin{definition}[CNN как функция в постановке Бьетти--Майраль]
Рассмотрим сигнал \(z_0 = x \in L^2(\mathbb{R}^d, \mathbb{R}^{p_0})\). Многоуровневое представление строится как:
\[
z_k = A_k \tilde{z}_k, \quad \tilde{z}_k(u) = \varphi_k(P_k z_{k-1}(u)), \quad k = 1, \ldots, n,
\]
где
\begin{itemize}
    \item \(P_k\) извлекает локальный патч вокруг позиции \(u\),
    \item \(\varphi_k\) --- \emph{нелинейность на патче}, индуцированная ядром скалярного произведения,
    \item \(A_k\) --- линейный оператор пулинга/усреднения.
\end{itemize}
Предиктор определяется линейным выходным слоем: \(f_\sigma(x) = \langle w_{n+1}, z_n \rangle\).
\end{definition}

\textbf{Ключевое предположение:} Нелинейность \(\varphi_k\) может быть записана в \emph{однородной форме}. Для патча \(z\):
\[
\varphi_k(z) = \|z\| \, \sigma\left( \frac{\langle g, z \rangle}{\|z\|} \right),
\]
где \(\sigma\) --- \emph{гладкая положительно однородная функция активации} (например, сглаженный вариант ReLU), удовлетворяющая условию \(\sigma(\alpha t) = \alpha \sigma(t)\) для \(\alpha \geq 0\).

\begin{theorem}[Вложение CNN в RKHS \cite{Bietti_Group}]
При выполнении условий согласованности между разложением ядра скалярного произведения и функцией \(\sigma\), CNN-предиктор \(f_\sigma\) принадлежит RKHS \(\mathcal{H}_{K_n}\), порождённому глубинным ядром \(K_n\).
\end{theorem}
\textbf{Вывод:} Это означает, что для широкого класса CNN с гладкими однородными активациями \emph{можно применять весь аналитический аппарат RKHS} (оценки устойчивости, теоремы об обобщении, основанные на норме) для их анализа.

Ядерная точка зрения Бьетти--Майраль позволяет получить \emph{априорные, не зависящие от данных} гарантии для представления \cite{Bietti_Group}:

\begin{enumerate}
    \item \textbf{Гарантии инвариантности и устойчивости к деформациям.} Для диффеоморфизмов \(L_\tau\) (деформаций, включая сдвиги) доказываются оценки в стиле Мала:
    \[
    \|\Phi_n(L_\tau x) - \Phi_n(x)\| \leq (C_1\|\nabla\tau\|_\infty + C_2\|\tau\|_\infty) \|x\|.
    \]
    Сдвигам соответствует \(\nabla\tau = 0\), что приводит к инвариантности, контролируемой членом \(C_2\|\tau\|_\infty\).

    \item \textbf{Сохранение информации.} Формулируются условия, при которых \emph{субдискретизация} (downsampling) в операторах \(A_k\) не приводит к потере информации сигнала (инъективность представления по модулю группы инвариантности).

    \item \textbf{Связь сложности, устойчивости и обобщения.} Если итоговый предиктор линеен в RKHS (\(f(x) = \langle w, \Phi_n(x) \rangle_{\mathcal{H}_{K_n}}\)), то, используя свойство воспроизведения и неравенство Коши-Буняковского:
    \[
    |f(x) - f(x')| \leq \|f\|_{\mathcal{H}_{K_n}} \|\Phi_n(x) - \Phi_n(x')\|_{\mathcal{H}_{K_n}}.
    \]
    Таким образом, \emph{устойчивость предсказаний наследует устойчивость представления}, а норма \(\|f\|_{\mathcal{H}_{K_n}}\) служит одновременно мерой сложности модели (регуляризатором) и контролем чувствительности.
\end{enumerate}

\subsection{Инвариантность и устойчивость в ядерном представлении}
Следуещим делом нам необъодимо подробно разобраться в механизмах, с помощью которых свёрточные представления, смоделированные в рамках подхода Бьетти--Майраль, обеспечивают инвариантность к действию групп (например, сдвигов) и устойчивость к малым деформациям сигнала. Ключевую роль в этих механизмах играют операции \emph{пулинга (усреднения)}.

Пусть группа \( G \) действует на пространстве сигналов \(\mathcal{X}\) с помощью операторов \( T_g \).
\begin{definition}[Групповой пулинг (усреднение по орбите)]
\label{def:group_pooling}
Для компактной группы \( G \) с вероятностной мерой Хаара \( \mu \) определим оператор \emph{группового пулинга} \( A_G: \mathcal{H} \to \mathcal{H} \):
\[
A_G(u) \coloneq \int_G \rho(g) u \, d\mu(g).
\]
Оператор \( A_G \) является \emph{линейным} и \emph{в точности инвариантным}:
\[
A_G(\rho(h) u) = A_G(u) \quad \forall h \in G, \forall u \in \mathcal{H}.
\]
\end{definition}
\textbf{Ключевое следствие:} Если признаковое отображение \( \Phi \) является \emph{эквивариантным}, то композиция \( x \mapsto A_G(\Phi(x)) \) будет \emph{инвариантной} функцией. Таким образом, в рамках ядерного подхода \cite{Bietti_Group} \emph{инвариантность возникает как результат применения линейного оператора усреднения \( A_G \) к эквивариантному представлению}.

\begin{example}[Глобальный пулинг для сдвигов]
Рассмотрим сигнал на решётке \( z\colon \mathbb{Z}^\nu \to \mathbb{R}^m \). Сдвиг на \( \tau \) задаётся оператором \( (R_\tau z)(\gamma) = z(\gamma - \tau) \). Простейший \emph{глобальный пулинг} --- это суммирование или усреднение по всему домену \( \Omega \):
\[
A_{\text{global}}(z) \coloneq \sum_{\gamma \in \Omega} z(\gamma) \quad \text{или} \quad \frac{1}{|\Omega|} \sum_{\gamma \in \Omega} z(\gamma).
\]
Если представление \( z \) является \emph{сдвиго-эквивариантным}, то \( A_{\text{global}}(z) \) будет \emph{сдвиго-инвариантным} на выбранном домене \( \Omega \).   
\end{example}

В реальных CNN пулинг является \emph{локальным} (усреднение или максимум по окну \( B \subset \mathbb{Z}^\nu \)) и часто сопровождается \emph{прореживанием (stride)} \( s \):
\[
(\text{Pool } z)(\gamma) \coloneq \text{pool}\{z(\gamma + \beta) : \beta \in B\}, \quad (\text{Down}_s z)(\gamma) \coloneq z(s\gamma).
\]

\begin{proposition}[Влияние пулинга и прореживания на эквивариантность]
Чистые свёртки сохраняют эквивариантность относительно сдвигов на \emph{исходной решётке}. После операции прореживания с шагом \( s \) эквивариантность сохраняется, но относительно \emph{более coarse решётки}: точными симметриями остаются только сдвиги, кратные \( s \) (подгруппа \( s\mathbb{Z}^\nu \)).
\end{proposition}
\textbf{Интуиция:} Локальный пулинг даёт \emph{приближённую (локальную) инвариантность}: малые сдвиги входа (меньшие, чем масштаб окна пулинга) приводят к малым изменениям в признаках, так как усреднение сглаживает высокочастотные компоненты. Точная инвариантность обычно требует \emph{глобального агрегирования} (суммирования по всей пространственной области или орбите группы).

Для анализа устойчивости в работе \cite{Bietti_Group} вводится модель малой деформации сигнала.

\begin{definition}[Оператор деформации]
\label{def:deformation_operator}
Пусть \( x: \mathbb{R}^d \to \mathbb{R}^p \) --- сигнал (изображение), а \( \tau: \mathbb{R}^d \to \mathbb{R}^d \) --- \emph{поле малых смещений} (деформация). \emph{Деформированный сигнал} определяется как:
\[
(L_\tau x)(u) \coloneq x(u - \tau(u)) = x \circ (\mathrm{Id} - \tau)(u).
\]
Мерой ``величины'' деформации служит норма:
\[
\| \tau \|_{\text{size}} \coloneq \|\tau\|_\infty + \|\nabla \tau\|_\infty,
\]
где \( \|\tau\|_\infty = \sup_u \|\tau(u)\| \), а \( \|\nabla \tau\|_\infty = \sup_u \|\nabla \tau(u)\|_{\text{op}} \). Сдвиги являются частным случаем при \( \tau(u) \equiv a \).
\end{definition}

В рамках многоуровневой конструкции представления \( \Phi_n \) (см. Раздел \ref{sec:5.2_conv_kernel_view}) целью является получение оценки, связывающей изменение представления с величиной деформации.

\begin{definition}[Липшицевость относительно деформаций]
\label{def:desired_stability_bound}
Для глубинного представления \( \Phi_n \) и оператора деформации \( L_\tau \) желательно получить оценку вида:
\[
\| \Phi_n(x) - \Phi_n(L_\tau x) \|_{L^2(\mathbb{R}^d; \mathcal{H}_n)} \leq C \| \tau \|_{\text{size}} \| x \|_{L^2(\mathbb{R}^d)},
\]
где константа \( C \) может зависеть от параметров сети (глубины \( n \), ядер пулинга \( h_k \)), но не от конкретных \( x \) и \( \tau \).
\end{definition}
\textbf{Интерпретация:} Малая деформация \( \tau \) (малая как по величине смещения \( \|\tau\|_\infty \), так и по его градиенту \( \|\nabla \tau\|_\infty \)) приводит к малому изменению глубинного представления. Это свойство \emph{устойчивости} критически важно для надёжности модели на слегка искажённых данных (например, слегка повёрнутых или деформированных изображениях).

\textbf{Роль операторов \( A_k \):} В конструкции Бьетти--Майраль \cite{Bietti_Group} операторы \( A_k \) (усреднения/пулинга) являются \emph{линейными} и, как правило, \emph{неэкспансивными} (сжимающими). Это свойство позволяет провести анализ, показывая, что сжатие информации на каждом уровне через усреднение подавляет чувствительность к высокочастотным искажениям, таким как мелкомасштабные деформации.

\subsection{Связь спектральной нормы и регуляризации}
Данный раздел завершает связь между ядерным представлением свёрточных сетей и практическими методами их обучения. Мы покажем, как норма в RKHS, возникающая в конструкции Бьетти--Майраль \cite{Bietti_Group}, естественным образом приводит к \emph{регуляризации спектральных норм весовых матриц} --- широко используемой эвристике в глубоком обучении для обеспечения устойчивости и обобщения.

В рамках ядерного подхода \cite{Bietti_Group} сложность CNN-предиктора естественно измеряется нормой в соответствующем RKHS.
\begin{proposition}[Ограничение нормы CNN в RKHS \cite{Bietti_Group}]
Пусть CNN-предиктор \( f_\sigma \), построенный с гладкой однородной функцией активации \( \sigma \), принадлежит RKHS \( \mathcal{H}_{K_n} \) глубинного ядра \( K_n \). Тогда существуют константы \( C_\sigma \), зависящие от активации и ядра, такие что выполняется оценка вида:
\begin{equation}
\label{eq:rkhs_norm_bound_cnn}
\| f_\sigma \|^2_{\mathcal{H}_{K_n}} \leq \| w_{n+1} \|_2^2 \, C_\sigma^2\!\left( \| W_n \|_2^2 \, C_\sigma^2\!\left( \| W_{n-1} \|_2^2 \, \cdots \, C_\sigma^2\!\left( \| W_1 \|_F^2 \right) \cdots \right) \right),
\end{equation}
где
\begin{itemize}
    \item \( w_{n+1} \) --- веса выходного (линейного) слоя,
    \item \( W_k \) --- матрица, кодирующая свёрточные фильтры \( k \)-го слоя,
    \item \( \| \cdot \|_2 \) --- спектральная (операторная) норма,
    \item \( \| \cdot \|_F \) --- норма Фробениуса.
\end{itemize}
\end{proposition}

\textbf{Интерпретация:}
\begin{itemize}
    \item Норма предиктора в RKHS ограничена \emph{вложенным произведением спектральных норм весовых матриц} всех слоёв (с коррекцией \( C_\sigma \) на каждом шаге).
    \item В \emph{линейном случае} (\( C_\sigma^2(\lambda^2) = \lambda^2 \)) оценка \eqref{eq:rkhs_norm_bound_cnn} упрощается до:
    \[
    \| f_\sigma \|_{\mathcal{H}_{K_n}} \lesssim \| w_{n+1} \|_2 \cdot \| W_n \|_2 \cdot \| W_{n-1} \|_2 \cdots \| W_1 \|_F.
    \]
    \item Таким образом, контролируя спектральные нормы \( \| W_k \|_2 \) и норму Фробениуса \( \| W_1 \|_F \), мы контролируем и норму в RKHS \( \| f_\sigma \|_{\mathcal{H}_{K_n}} \), которая является естественной мерой сложности модели в рамках ядерной теории.
\end{itemize}

Ограничение спектральных норм весовых матриц --- распространённая практика в глубоком обучении. Ядерный подход Бьетти--Майраль даёт этому две взаимодополняющие теоретические интерпретации.\\
\textbf{(A) Управление устойчивостью через контроль числа Липшица.}
Для композиции функций \( f = F_n \circ \cdots \circ F_1(x) \) (например, слоёв CNN с 1-липшицевыми активациями типа ReLU) выполняется:
\[
\mathrm{Lip}(f) \leq \prod_{k=1}^{n} \mathrm{Lip}(F_k) \lesssim \prod_{k=1}^{n} \| W_k \|_2.
\]
Следовательно, ограничение произведения спектральных норм \( \prod_{k=1}^{n} \| W_k \|_2 \) напрямую ограничивает \emph{константу Липшица всей сети}:
\[
\| f(x) - f(x') \| \leq \left( \prod_{k=1}^{n} \| W_k \|_2 \right) \| x - x' \|.
\]
\textbf{Практический смысл:} Малые возмущения или деформации входа приводят к ограниченным изменениям выхода, что повышает \emph{устойчивость (robustness)} модели.

\noindent
\textbf{(B) Управление сложностью через норму в RKHS (ядро Бьетти--Майраль).}
Как показано в оценке \eqref{eq:rkhs_norm_bound_cnn}, нормы \( \| W_k \|_2 \) (и \( \| W_1 \|_F \)) входят в верхнюю границу для нормы предиктора в RKHS \( \mathcal{H}_{K_n} \). В задачах обучения с регуляризацией минимизируется функционал вида:
\[
\frac{1}{n} \sum_{i=1}^{n} \ell(y_i, f(x_i)) + \lambda \| f \|^2_{\mathcal{H}_{K_n}}.
\]
Таким образом, ограничение спектральных норм является \emph{естественным прокси} для ограничения RKHS-нормы, которая, в свою очередь, контролирует сложность модели и способствует \emph{обобщению (generalization)}.

Из всего вышесказанного можно сделать вывод: регуляризация спектральных норм --- это не просто эвристика, а метод, имеющий строгое теоретическое обоснование в рамках ядерного подхода к CNN.
\begin{itemize}
    \item \textbf{Способствует устойчивости:} Ограничивает расширение признаков (константу Липшица сети), делая предсказания менее чувствительными к малым возмущениям входных данных.
    \item \textbf{Согласуется с контролем сложности в RKHS:} Является практической реализацией ограничения нормы в гильбертовом пространстве, порождённом свёрточным ядром.
    \item \textbf{Поддерживает обобщающую способность:} Прямо следует из теории обобщения для kernel methods, где норма в RKHS фигурирует в классических оценках типа Радемахера.
\end{itemize}

\textbf{Итог:} В работе Бьетти--Майраль \cite{Bietti_Group} устанавливается мост между \emph{геометрическими свойствами представления} (инвариантность, устойчивость) и \emph{сложностью модели} (регуляризация). Норма в RKHS становится единой мерой, связывающей эти аспекты: её ограничение одновременно способствует устойчивости представления (через оценку Липшица) и улучшает обобщение (через контроль capacity модели).

\subsection{Информационная сохранность и субдискретизация}

Теперь рассмотрим критический аспект проектирования многоуровневых представлений: компромисс между достижением инвариантности (через пулинг и субдискретизацию) и сохранением полезной информации из входного сигнала. В рамках подхода Бьетти--Майраль \cite{Bietti_Group} анализируются условия, при которых субдискретизация не приводит к катастрофической потере информации.

В CNN операции пулинга (усреднения) и субдискретизации (downsampling) служат двум основным целям:
\begin{enumerate}
    \item Увеличить \emph{инвариантность} представления к малым пространственным сдвигам.
    \item Уменьшить пространственное разрешение и, следовательно, вычислительную сложность последующих слоёв.
\end{enumerate}

Однако существует фундаментальный \emph{компромисс}:
\begin{itemize}
    \item \textbf{Агрессивный пулинг / сильная субдискретизация:} Обеспечивает высокую степень инвариантности, но может привести к \emph{неинъективности (aliasing)} --- различным входным сигналам будет соответствовать одно и то же представление, то есть информация будет безвозвратно потеряна.
    \item \textbf{Слабый пулинг / сохранение разрешения:} Сохраняет детальную пространственную информацию, но обеспечивает меньшую инвариантность и вычислительную эффективность.
\end{itemize}
Ключевой вопрос, исследуемый в \cite{Bietti_Group}: \emph{Можно ли добиться одновременно устойчивости (инвариантности) к преобразованиям и сохранения информации?}

Напомним общую схему построения представления \cite{Bietti_Group}:
\[
x_0 \coloneq x, \quad x_k \coloneq A_k M_k P_k x_{k-1}, \quad k = 1, \ldots, n,
\]
где \(P_k\) извлекает патчи, \(M_k\) --- ядерная нелинейность, а \(A_k\) --- оператор усреднения (пулинга), за которым часто следует \emph{субдискретизация (downsampling)}.

\textbf{Декомпозиция \(A_k\):} Оператор \(A_k\) можно представить как композицию \(A_k \approx D_k \circ S_k\), где:
\begin{itemize}
    \item \(S_k\) --- \emph{локальное усреднение (сглаживание)} с ядром \(h_k\) (низкочастотный фильтр, антиалиасинг).
    \item \(D_k\) --- собственно \emph{операция субдискретизации} (децимация) с заданным шагом (stride).
\end{itemize}

Когда субдискретизация приводит к потере информации?
Она теряется безвозвратно (происходит \emph{алиасинг}), когда:
\begin{enumerate}
    \item \textbf{Нарушение условия Найквиста:} Высокочастотные компоненты в признаковой карте \(M_k P_k x_{k-1}\) имеют частоту выше \emph{половины частоты дискретизации}, задаваемой шагом \(D_k\).
    \item \textbf{Слишком большой шаг (stride):} Шаг субдискретизации значительно превышает \emph{корреляционную длину} признаков (масштаб, на котором признаки существенно изменяются).
    \item \textbf{Отсутствие избыточности:} Патчи, извлекаемые оператором \(P_k\), почти не перекрываются, и потерянные при субдискретизации координаты не могут быть восстановлены из соседних.
\end{enumerate}

Субдискретизация не приводит к существенной потере информации (с точностью до предполагаемых инвариантностей), если выполняются следующие условия, вытекающие из анализа в \cite{Bietti_Group}:
\begin{enumerate}
    \item \textbf{Предварительное сглаживание (Антиалиасинг):} Оператор \(S_k\) (часть \(A_k\)) должен эффективно \emph{ограничивать полосу частот} признаковой карты \(M_k P_k x_{k-1}\), делая её \emph{приблизительно ограниченной по спектру (band-limited)} на масштабе, соответствующем шагу дискретизации \(D_k\).
    \item \textbf{Согласование шага и полосы:} Шаг субдискретизации \(D_k\) должен быть \emph{согласован с эффективной шириной полосы} сглаженного сигнала (интуиция Найквиста-Шеннона).
    \item \textbf{Избыточность через перекрытие патчей:} Оператор \(P_k\) должен извлекать патчи с \emph{большим перекрытием}, обеспечивая избыточность в представлении, которая позволяет восстановить информацию, усреднённую на этапе \(S_k\).
    \item \textbf{Выразительность признакового отображения:} Ядерное отображение \(M_k\) (и всё последующее представление) должно быть достаточно выразительным, чтобы \emph{различать различные локальные патчи}, даже после усреднения.
\end{enumerate}

В работе Бьетти--Майраль <<сохранение информации>> формализуется через свойство \emph{инъективности по модулю группы инвариантности}.
\begin{definition}[Инъективность по модулю группы]
\label{def:injective_modulo_group}
Пусть \(G\) --- группа преобразований (например, сдвигов). Представление \(\Phi\) называется \emph{инъективным по модулю \(G\)}, если для любых входных сигналов \(x, x'\) выполняется:
\[
\Phi(x) = \Phi(x') \quad \text{тогда и только тогда, когда} \quad x' = T_g x \text{ для некоторого } g \in G.
\]
Эквивалентно: \(\|\Phi(x) - \Phi(x')\|\) мало тогда и только тогда, когда \(x\) и \(x'\) близки с точностью до действия группы \(G\).
\end{definition}

\textbf{Итог:} В успешно спроектированной многоуровневой архитектуре, следующей принципам \cite{Bietti_Group}, операция \(A_k = D_k \circ S_k\) работает как \emph{координатное преобразование}, которое:
\begin{enumerate}[label=\arabic*)]
    \item \emph{Сглаживает} данные для подавления чувствительности к высокочастотным искажениям (деформациям),
    \item \emph{Дискретизирует} с подходящей частотой, сохраняя низкочастотную (смысловую) информацию,
    \item В идеале сохраняет свойство \emph{инъективности по модулю группы} инвариантности, то есть разные входы (не эквивалентные относительно \(G\)) отображаются в разные представления.
\end{enumerate}
Таким образом, пулинг и субдискретизация в такой модели --- не просто механическое уменьшение размерности, а часть стратегии построения \emph{устойчивого, инвариантного, но при этом информативного} представления сигнала.

\end{document}